\section{Proofs for the Strongly Convex Case (Section \ref{sec:strongly-convex})}
\label{app:strongly-convex}

\begin{lemma}
\label{lem:basic-lemma-sc}
For any sequence $\weight_t>0, g_t\in\RR^d$, suppose an online learner predicts $w_t$ and receives $t$-th loss $\ell_t(w)=\langle g_t,w\rangle$, and define $x_t=\sum_{i=1}^t\frac{\weight_iw_i}{\weight_{1:t}}$. If $f$ is $\mu$-strongly convex w.r.t. $\|\cdot\|$, then
\begin{align*}
    \weight_{1:T}(f(x_T) - f(x^*)) \leq \regret_T(x^*) + \sum_{t=1}^T \left( \langle \weight_t \nabla f(x_t) - g_t, w_t-x^* \rangle - \frac{\weight_t\mu}{2}\|x_t-x^*\|^2 \right).
\end{align*}
\end{lemma}

\begin{proof}
We start with the strong convexity identity $f(x^*)\geq f(x_t)+\langle\nabla f(x_t),x^*-x_t\rangle + \frac{\mu}{2}\|x_t-x^*\|^2$:
\begin{align}
    \sum_{t=1}^T\weight_t(f(x_t) - f(x^*))
    &\leq \sum_{t=1}^T \weight_t\langle \nabla f(x_t),x_t-x^*\rangle - \frac{\weight_t\mu}{2}\|x_t-x^*\|^2.
    \label{eq:sc-1}
\end{align}
With the same argument in the proof of Lemma \ref{lem:basic-lemma}, we can show:
\begin{align*}
    \weight_t\langle \nabla f(x_t),x_t-x^*\rangle
    &= \weight_t \langle \nabla f(x_t), x_t-w_t \rangle + \weight_t \langle \nabla f(x_t), w_t-x^*\rangle 
    \intertext{Recall the definition that $\weight_{1:t}x_t = \weight_{1:t-1}x_{t-1}+\weight_tw_t$ and thus $\weight_t(x_t-w_t)=\weight_{1:t-1}(x_{t-1}-x_t)$. Also, since $f$ is convex, $\langle \nabla f(x_t),x_{t-1}-x_t\rangle \leq f(x_{t-1})-f(x_t)$, so:}
    &\leq \weight_{1:t-1}f(x_{t-1}) - \weight_{1:t-1}f(x_t) + \langle \weight_t\nabla f(x_t) - g_t + g_t,w_t-x^*\rangle.
\end{align*}
Consequently, moving $\sum_{t=1}^T \weight_tf(x_t)$ to the right side and taking the telescopic sum in \eqref{eq:sc-1} gives:
\begin{align*}
    -\weight_{1:T}f(x^*)
    &\leq \sum_{t=1}^T \weight_t\langle \nabla f(x_t),x_t-x^*\rangle -\weight_tf(x_t) - \frac{\weight_t\mu}{2}\|x_t-x^*\|^2 \\
    &\leq -\weight_{1:T}f(x_T) + \regret_T(x^*) + \sum_{t=1}^T \langle\weight_t\nabla f(x_t)-g_t,w_t-x^*\rangle - \frac{\weight_t\mu}{2}\|x_t-x^*\|^2.
\end{align*}
Moving $\weight_{1:T}f(x_T)$ to the left completes the proof.
\end{proof}

This lemma immediately implies Lemma \ref{lem:strongly-convex-regularized-loss}.

\begin{lemma}
\label{lem:DP-OTB:strongly-convex-regularized-loss}
Suppose $f$ is $\mu$-strongly convex w.r.t. $\|\cdot\|$. If we replace $\ell_t(w)=\langle g_t+\gamma_t,w\rangle$ with $\bar \ell_t(w)=\ell_t(w)+\frac{\weight_t\mu}{4}\|w-x_t\|^2$ in Algorithm \ref{alg:DP-OTB:main}, and denote the associated regret by $\overline \regret_T$, then
\begin{align*}
    \weight_{1:T} ( f(x_T) -  f(x^*))
    &\leq \overline\regret _T(x^*) +\sum_{t=1}^T \langle \weight_t\nabla f(x_t)-g_t-\gamma_t, w_t-x^*\rangle - \frac{\weight_t\mu}{8}\|w_t-x^*\|^2 \\
    &\leq \overline\regret_T(x^*) + \sum_{t=1}^T \frac{2\|\weight_t\nabla f(x_t)-g_t-\gamma_t\|_*^2}{\weight_t \mu}.
\end{align*}
\end{lemma}

\begin{proof}
By definition, $\bar{\ell}_t(w)=\ell_t(w) + \frac{\weight_t\mu}{4}\|w-x_t\|^2$, so
\begin{align*}
    \overline \regret_T(x^*) 
    &= \sum_{t=1}^T \left( \ell_t(w_t)+\frac{\weight_t\mu}{4}\|w_t-x_t\|^2\right) - \left( \ell_t(x^*) + \frac{\weight_t\mu}{4}\|x^*-x_t\|^2 \right) \\
    &= \regret_T(x^*) + \sum_{t=1}^T \frac{\weight_t\mu}{4}(\|w_t-x_t\|^2-\|x_t-x^*\|^2).
\end{align*}
Upon substituting this equation into Lemma \ref{lem:basic-lemma-sc}, we get:
\begin{align*}
    &\quad \ \weight_{1:T}(f(x_T) - f(x^*)) \\
    &\leq \overline\regret_T(x^*) -\sum_{t=1}^T \frac{\weight_t\mu}{4}(\|w_t-x_t\|^2-\|x_t-x^*\|^2) \\
    &+ \sum_{t=1}^T \langle \weight_t \nabla f(x_t)-g_t-\gamma_t, w_t-x^*\rangle - \frac{\weight_t\mu}{2} \|x_t-x^*\|^2 \\
    &\leq \overline\regret_T(x^*) +\sum_{t=1}^T \langle \weight_t\nabla f(x_t)-g_t-\gamma_t, w_t-x_t\rangle - \frac{\weight_t\mu}{4}(\|w_t-x_t\|^2+\|x_t-x^*\|^2) \\
    &\leq \overline\regret_T(x^*) +\sum_{t=1}^T \langle\weight_t\nabla f(x_t)-g_t-\gamma_t, w_t-x^*\rangle - \frac{\weight_t\mu}{8}\|w_t-x^*\|^2.
\end{align*}
The last inequality follows from the identity $\|w_t-x^*\|^2 \leq 2\|w_t-x_t\|^2 + 2\|x_t-x^*\|^2$.

For the second inequality in the lemma, by Fenchel-Young's inequality,
\begin{align*}
    &\langle\weight_t\nabla f(x_t)-g_t-\gamma_t, w_t-x^*\rangle - \frac{\weight_t\mu}{8}\|w_t-x^*\|^2 \\
    \leq &\|\weight_t\nabla f(x_t)-g_t-\gamma_t\|_*\|w_t-x^*\|-\frac{\weight_t\mu}{8}\|w_t-x^*\|^2 \\
    % \intertext{Denote $B_t=\|\weight_t\nabla f(x_t)-g_t-\gamma_t\|_*$ for short, and note that this is a quadratic equation:}
    \intertext{For any quadratic of form $ax-bx^2$ and $a,b>0$, note that $\sup_x ax-bx^2\leq a^2/4b$, so:}
    \leq & \frac{2\|\weight_t\nabla f(x_t)-g_t-\gamma_t\|_*^2}{\weight_t\mu}.
\end{align*}
\end{proof}

\begin{proposition}
\label{prop:subdiff-chain-rule}
Suppose $W$ is a convex bounded domain with diameter $D$, and let $u\in W$ and $f(w)=\|w-u\|^2$. Then for all $w\in W$ and $v\in \partial f(w)$, $\|v\|_*\leq 2D$.
\end{proposition}

\begin{proof}
Let $\phi(r)=r^2$ and $g(w)=\|w-u\|$, then $f(w) = \phi\circ g(w)$. By chain rule of sub-differentials (Corollary 16.72 \cite{bauschke2011convex}), 
\begin{align*}
    \partial f(w)
    &= \{\alpha v' : \alpha\in \partial \phi(g(w)), v'\in\partial g(w) \} \\
    &= \{2\|w-u\| v' : v'\in\partial \|w-u\|\}.
\end{align*}
By assumption, $\|w-u\|\leq D$. Moreover, $\|\cdot\|$ is $1$-Lipschitz (because $\|x\|-\|y\|\leq \|x-y\|$), so $\|v'\|_*\leq 1$ for all $v'\in\partial \|w-u\|$. As a result, for all $v\in \partial f(w)$, $\|v\|_* = 2\|w-u\|\|v'\|_* \leq 2D$.
\end{proof}